\scanpage{77}
Ainsi
\[
 0\leq\alpha\bar\alpha\leq1
\]
et de façon précise
\[
 1-\alpha\bar\alpha=4\rho\sigma,\quad\text{i.e.}\quad \alpha\bar\alpha=1-4\rho\sigma
\]

Ainsi, $\beta$ et $\alpha\bar\alpha$ sont donnés en termes volumétriques. Il reste à donner (pour $\alpha\bar\alpha\unsure{\ne|+}0$) $\arg\alpha$ (qui dépend bien \unsure{sur|sûr} encore du choix de $\omega$). Choisissant $\omega$ de façon que $\omega^{12}=\Delta$, ce qui \unsure{le|la} détermine mod multiplication par $\zeta$, $\zeta^{12}=1$, on voit que $\arg\alpha\in U=\{z\in\mathbb C,\lvert z\rvert=1\}$ se trouve déterminé mod \deleted{\unsure{puisse|puissance}} mult. par $\unsure{\xi|\zeta}$, avec $\xi^6=1$.

On a donc comme invariant canonique, (déterminé \unsure{directement|comment} en fonction de $j$ !) $(\arg\alpha)^6\in U$. NB. $\rho$ et $\sigma$ sont déterminés mod rien du tout par la normalisation $\omega^{12}=\Delta$).
% archive sheet number: 76
